This work studies the reliability of autonomous web agents and therefore touches
on the risks of deploying them.

\emph{Automated browsing.} All experiments ran against self-hosted WebArena
replicas with synthetic data on isolated infrastructure. No live third-party site
was traversed, no real account was used, and no external service was affected.
The authentication state used by the agent exists only inside that sandbox.

\emph{Fault injection.} The injected faults --- dropped DOM nodes, blocking
overlays, HTTP errors, timeouts, and corrupted tool arguments --- are
non-adversarial: they degrade the agent's own inputs and do not target, probe, or
exfiltrate anything. The injection harness is deterministic and seeded, and it is
released so that its effects can be reproduced rather than speculated about.

\emph{Dual use.} A systematic account of which faults a web agent fails to notice
could in principle inform an attacker seeking to induce failures, or a site
operator seeking to block agents. We judge the risk low relative to the benefit:
the faults are already ubiquitous in ordinary web operation, the findings are
mostly about measurement rather than about exploitable weaknesses, and the
mitigation we argue for --- process-level monitoring and evaluator hygiene ---
helps defenders. We release no attack tooling and no adversarial payloads.

\emph{Evaluation integrity.} The paper's central methodological claim is that a
widely used evaluation setup can be wrong in a way that flatters the system being
measured, and that a structurally failure-biased evaluator path went undetected.
We report this because it affects how published numbers should be read, not to
single out a particular benchmark; the failure is in the reporting convention, not
in the benchmark's design.
